\section{Introduction}
\label{sec:intro}

Software evolves continuously.  Regression testing is therefore not a secondary
engineering convenience: it is the mechanism by which evidence about an earlier
version is reused to decide whether a modified version still satisfies its
requirements.  Classical regression-test selection and change-impact analysis
make this reuse explicit \cite{RothermelHarrold1996,YooHarman2012,LiSunLeungZhang2013},
and recent incremental-verification work similarly reuses proof artifacts across
program versions \cite{YinEtAl2025}.

Fairness creates a regression requirement with an additional statistical
challenge.  Software fairness testing is established as a testing problem
\cite{GalhotraBrunMeliou2017}, while online fairness auditing studies how observations can be adaptively allocated
to assess a fixed deployed system \cite{ManerikerBurleyParthasarathy2023}.  Yet
software may itself be nondeterministic even in the complete absence of learned
models.  Randomized tie-breaking, task scheduling, concurrent interleavings,
resource contention, random backoff, clocks, and other nondeterministic APIs can
make repeated executions of the same code and input disagree.  Recent testing
work continues to find nondeterminism central to flakiness \cite{YuanLinShi2026}.

We ask a version-aware question:
\begin{quote}
\emph{Given a previously certified program, a patch, and stochastic execution
access to the old and new versions, what is the minimum number of executions
required to detect or rule out a fairness regression?}
\end{quote}

The regression setting contributes two pieces of information absent from fresh
auditing.  First, a sound change-impact analysis can localize where behavior may
have changed.  Second, the old version remains executable.  If stochastic choices
can be replayed through a controlled seed or random tape, $P_0$ and $P_1$ can be
run under the same execution randomness.  Their difference then estimates the
patch itself.  This is related to the classical common-random-numbers principle in
simulation \cite{GlassermanYao1992}, and to coupling as a formal technique for
reasoning about pairs of probabilistic programs \cite{BartheEtAl2017}; our use is
as an information model for regression testing.

We represent fairness requirements by a counterfactual graph $G=(V,E)$.  A vertex
is a valid software input and an edge is a pair whose expected outcomes should
remain sufficiently close.  A patch produces expected behavioral changes
$\theta_v=\mu_{1,v}-\mu_{0,v}$.  The central core problem asks whether
$\theta_u=\theta_v$ on every impacted fairness edge or whether at least one edge
has differential patch effect at least $\gamma$.  This \emph{patch-neutrality}
problem is stronger than full thresholded fairness preservation but isolates the
information-theoretic interaction among version reuse, graph structure, and
stochastic execution.

For Gaussian paired deltas we derive a graph complexity
\[
\Psi(H,\nu)=
\min_{q_v>0}\left\{
\sum_v q_v:
\frac{\nu_u^2}{q_u}+\frac{\nu_v^2}{q_v}\le1
\ \forall (u,v)\in E(H)
\right\},
\]
where $\nu_v^2$ is the variance of the old--new behavioral difference at vertex
$v$.  A static allocation based on $\Psi$ gives the upper bound, while an adaptive
change-of-measure argument gives a matching dependence on $\Psi$ and
$\gamma^{-2}$ in the lower bound.  We also derive the Lagrange dual, making clear
that the same executions can support several incident counterfactual comparisons.

The regression advantage appears directly in the variance.  If $Y_0,Y_1$ are old
and new outcomes produced under the same random tape,
\[
\mathrm{Var}(Y_1-Y_0)
=
\mathrm{Var}(Y_1)+\mathrm{Var}(Y_0)-2\mathrm{Cov}(Y_1,Y_0).
\]
When nearby versions respond similarly to the same stochastic environment, the
covariance can be large and the patch can be substantially easier to estimate
than either absolute program behavior.  In a common-uniform Bernoulli model, a
change from success probability $p$ to $p+d$ yields differential variance
$d(1-d)$, which vanishes with patch size even when the marginal outcome variances
remain bounded away from zero.


Our Bernoulli results expose a second form of reuse.  Four-way counterfactual
replay synchronizes randomness across the old and new versions and both endpoints
of one fairness edge.  A fairness-neutral common patch effect then cancels
pathwise, while a differential effect of size \(d\) remains a sparse witness;
this changes fixed-confidence complexity from quadratic to linear in \(1/d\).
If the test harness can replay one seed across the entire fairness-active input
suite, the same round simultaneously exposes every counterfactual edge.  For the
homogeneous family analyzed in \Cref{thm:suite-replay}, the number of replay
rounds is independent of \(|E|\), and concrete execution complexity scales with
the number of active inputs rather than the number of fairness comparisons.
This stronger oracle is stated explicitly: it is not assumed available to every
software system.

Between edgewise replay and full-suite replay lies a partial-replay regime in
which only declared groups of inputs can share a tape.  Here two combinatorial
objects diverge.  A four-way edge witness requires a replay group containing both
fairness endpoints, whereas raw transcript information requires only the changed
active endpoint plus a synchronized reference.  We show that the corresponding
edge-cover geometry can exceed the active-vertex cover $C_{\rm vert}$ by an
unbounded factor and can even be infeasible.  More importantly, for the frozen
homogeneous single-changed-endpoint Bernoulli family, the raw information is
operational: scheduling groups by $C_{\rm vert}$ and rejecting on a nonconstant
within-group vector of old--new patch deltas yields zero type-I error and
$O(C_{\rm vert}d^{-1}\log(1/\delta))$ fixed-design cost, while transcript
change of measure gives the matching $C_{\rm vert}$ combinatorial lower
geometry.  This theorem is deliberately not extended to heterogeneous
fairness-preserving nulls, where unrelated synchronized references can undergo
different benign patch effects.

The paper is deliberately theory-first and model-free.  No pretrained model,
open-source model, proprietary API, or LLM is an experimental subject.  GPU
computation is used only to falsify conjectures, validate implementations, and
stress-test the proved complexity laws on large synthetic instances.
